\chapter*{Presentación del curso}
\addcontentsline{toc}{chapter}{Presentación del curso}\markboth{Presentación del curso}{}

\textbf{Perfil:} capacitación en datos/analítica · \textbf{Nivel:}
principiante → intermedio · \textbf{Duración:} 8 módulos,
\textasciitilde34 h · \textbf{Modalidad:} autoformación o híbrida
(sesiones de 2 h + práctica) · \textbf{Idioma:} español

\section*{¿Para quién es?}
\addcontentsline{toc}{section}{¿Para quién es?}

Personas sin experiencia previa en programación (o que vienen de Excel)
que quieren usar R para limpiar, analizar, visualizar y reportar datos.
Los ejemplos usan un conjunto de \textbf{ventas simuladas de retail} (12
tiendas, 3 canales, 4 regiones, 5 productos, 52 semanas de 2025).

\section*{Objetivos del curso (diseño inverso)}
\addcontentsline{toc}{section}{Objetivos del curso (diseño inverso)}

Al terminar, la persona podrá:

\begin{enumerate}
\def\labelenumi{\arabic{enumi}.}
\tightlist
\item
  \textbf{Escribir} scripts de R legibles usando objetos, estructuras de
  datos, control de flujo y funciones propias.
\item
  \textbf{Importar y ordenar} datos de CSV/Excel siguiendo los
  principios de \emph{tidy data}.
\item
  \textbf{Transformar} y \textbf{resumir} datos con
  \passthrough{\lstinline!dplyr!} para calcular indicadores por grupo.
\item
  \textbf{Construir} visualizaciones claras con
  \passthrough{\lstinline!ggplot2!} adecuadas a cada pregunta.
\item
  \textbf{Aplicar e interpretar} estadística descriptiva, pruebas de
  hipótesis y regresión lineal.
\item
  \textbf{Producir} un reporte reproducible con Quarto que responda una
  pregunta de negocio.
\end{enumerate}

\textbf{Evidencia de logro:} ejercicios con solución en cada módulo, 3
entregas parciales (módulos 3, 5 y 7) y un \textbf{proyecto integrador}
evaluado con rúbrica (módulo 8).

\section*{Mapa del curso}
\addcontentsline{toc}{section}{Mapa del curso}

\begin{longtable}[]{@{}
  >{\raggedright\arraybackslash}p{(\columnwidth - 8\tabcolsep) * \real{0.2000}}
  >{\raggedright\arraybackslash}p{(\columnwidth - 8\tabcolsep) * \real{0.2000}}
  >{\raggedright\arraybackslash}p{(\columnwidth - 8\tabcolsep) * \real{0.2000}}
  >{\raggedright\arraybackslash}p{(\columnwidth - 8\tabcolsep) * \real{0.2000}}
  >{\raggedright\arraybackslash}p{(\columnwidth - 8\tabcolsep) * \real{0.2000}}@{}}
\toprule\noalign{}
\begin{minipage}[b]{\linewidth}\raggedright
\#
\end{minipage} & \begin{minipage}[b]{\linewidth}\raggedright
Módulo
\end{minipage} & \begin{minipage}[b]{\linewidth}\raggedright
Horas
\end{minipage} & \begin{minipage}[b]{\linewidth}\raggedright
Nivel de Bloom
\end{minipage} & \begin{minipage}[b]{\linewidth}\raggedright
Evaluación
\end{minipage} \\
\midrule\noalign{}
\endhead
\bottomrule\noalign{}
\endlastfoot
1 & Primeros pasos con R y RStudio & 4 & Recordar · Aplicar & Quiz \\
2 & Estructuras de datos & 4 & Comprender · Aplicar & Quiz + práctica \\
3 & Programación: control de flujo y funciones & 4 & Aplicar · Analizar
& \textbf{Entrega 1} (rúbrica) \\
4 & Importar y ordenar datos & 4 & Comprender · Aplicar & Quiz +
práctica \\
5 & Transformar datos con dplyr & 5 & Aplicar · Analizar &
\textbf{Entrega 2}: mini-reporte de KPIs \\
6 & Visualización con ggplot2 & 4 & Aplicar · Evaluar & Revisión entre
pares \\
7 & Estadística y modelos & 4 & Aplicar · Evaluar & \textbf{Entrega 3}:
informe breve \\
8 & Quarto y proyecto integrador & 5 & Crear & \textbf{Proyecto final}
(rúbrica global) \\
\end{longtable}

\section*{Ponderación sugerida}
\addcontentsline{toc}{section}{Ponderación sugerida}

\begin{longtable}[]{@{}ll@{}}
\toprule\noalign{}
Componente & Peso \\
\midrule\noalign{}
\endhead
\bottomrule\noalign{}
\endlastfoot
Quizzes y ejercicios (M1, M2, M4, M6) & 15 \% \\
Entrega 1 --- funciones (M3) & 15 \% \\
Entrega 2 --- KPIs con dplyr (M5) & 20 \% \\
Entrega 3 --- informe estadístico (M7) & 15 \% \\
Proyecto integrador (M8) & 35 \% \\
\end{longtable}

\section*{Preparación}
\addcontentsline{toc}{section}{Preparación}

\begin{enumerate}
\def\labelenumi{\arabic{enumi}.}
\item
  Instala \textbf{R} (\url{https://cran.r-project.org/}) y
  \textbf{RStudio Desktop}
  (\url{https://posit.co/download/rstudio-desktop/}).
\item
  Instala los paquetes del curso:

\begin{lstlisting}[language=R]
install.packages(c("tidyverse", "readxl", "broom", "scales", "renv"))
\end{lstlisting}
\item
  Abre esta carpeta como \textbf{Proyecto de RStudio} y genera los datos
  de práctica (ya incluidos en \passthrough{\lstinline!datos/!}):

\begin{lstlisting}[language=R]
source("datos/generar_ventas.R")
\end{lstlisting}
\end{enumerate}

Todo el código del curso se ejecutó y verificó con R 4.3.3 y tidyverse
2.0.0 usando este directorio como directorio de trabajo.

\section*{Estructura de archivos}
\addcontentsline{toc}{section}{Estructura de archivos}

\begin{lstlisting}
curso-r/
|-- README.md                          # este índice
|-- modulo-01-primeros-pasos.md
|-- modulo-02-estructuras-de-datos.md
|-- modulo-03-programacion.md
|-- modulo-04-importar-y-ordenar.md
|-- modulo-05-transformar-con-dplyr.md
|-- modulo-06-visualizacion-ggplot2.md
|-- modulo-07-estadistica-y-modelos.md
|-- modulo-08-reportes-quarto-y-proyecto.md
|-- datos/
|   |-- generar_ventas.R               # script reproducible (semilla fija)
|   |-- ventas.csv                     # 3 120 filas: fecha, tienda, producto, precio, unidades
|   `-- tiendas.csv                    # 12 tiendas: canal y región
`-- latex/
    |-- curso-r.tex                    # documento principal (portada, índice, preámbulo)
    |-- capitulos/                     # un .tex por módulo (generados)
    |-- convertir.py                   # regenera capitulos/ desde los .md (requiere pandoc)
    `-- curso-r.pdf                    # versión compilada
\end{lstlisting}

\textbf{Versión LaTeX/PDF:} desde \passthrough{\lstinline!latex/!},
ejecuta \passthrough{\lstinline!python3 convertir.py!} (solo si editaste
los \passthrough{\lstinline!.md!}) y luego
\passthrough{\lstinline!latexmk -pdf curso-r.tex!}. También se puede
subir la carpeta \passthrough{\lstinline!latex/!} a Overleaf y compilar
con pdfLaTeX.

\section*{Bibliografía general}
\addcontentsline{toc}{section}{Bibliografía general}

Leyenda: {[}ES{]} español · {[}EN{]} inglés · {[}OA{]} acceso abierto.

\subsection*{Libros}
\addcontentsline{toc}{subsection}{Libros}

\begin{itemize}
\tightlist
\item
  Grolemund, G. (2014). \emph{Hands-On Programming with R}. O'Reilly
  Media. ISBN 978-1-4493-5901-0.
  \url{https://rstudio-education.github.io/hopr/} {[}EN{]} {[}OA en
  línea{]}
\item
  Wickham, H. (2019). \emph{Advanced R} (2.ª ed.). CRC Press. ISBN
  978-0-8153-8457-1. \url{https://adv-r.hadley.nz/} {[}EN{]} {[}OA en
  línea{]}
\item
  Wickham, H., Çetinkaya-Rundel, M., \& Grolemund, G. (2023). \emph{R
  for Data Science: Import, tidy, transform, visualize, and model data}
  (2.ª ed.). O'Reilly Media. ISBN 978-1-4920-9740-2.
  \url{https://r4ds.hadley.nz/} {[}EN{]} {[}OA en línea{]}

  \begin{itemize}
  \tightlist
  \item
    Traducción al español (2.ª ed.):
    \url{https://davidrsch.github.io/r4ds-es/} {[}ES{]} {[}OA{]}
  \item
    Traducción comunitaria (1.ª ed.): \url{https://es.r4ds.hadley.nz/}
    {[}ES{]} {[}OA{]}
  \end{itemize}
\end{itemize}

\subsection*{Artículos}
\addcontentsline{toc}{subsection}{Artículos}

\begin{itemize}
\tightlist
\item
  Wickham, H. (2014). Tidy data. \emph{Journal of Statistical Software,
  59}(10), 1--23. \url{https://doi.org/10.18637/jss.v059.i10} {[}EN{]}
  {[}OA{]}
\item
  Wickham, H., Averick, M., Bryan, J., Chang, W., McGowan, L. D.,
  François, R., \ldots{} Yutani, H. (2019). Welcome to the tidyverse.
  \emph{Journal of Open Source Software, 4}(43), 1686.
  \url{https://doi.org/10.21105/joss.01686} {[}EN{]} {[}OA{]}
\end{itemize}

\subsection*{Documentación}
\addcontentsline{toc}{subsection}{Documentación}

\begin{itemize}
\tightlist
\item
  R Core Team. \emph{An Introduction to R}.
  \url{https://cran.r-project.org/doc/manuals/r-release/R-intro.html}
  {[}EN{]} {[}OA{]}
\item
  Quarto. \emph{Tutorial: Computations}.
  \url{https://quarto.org/docs/get-started/computations/rstudio}
  {[}EN{]}
\end{itemize}

\subsection*{Videos}
\addcontentsline{toc}{subsection}{Videos}

\begin{itemize}
\tightlist
\item
  freeCodeCamp.org. \emph{R Programming Tutorial -- Learn the Basics of
  Statistical Computing} {[}Video{]} (Barton Poulson, \textasciitilde2
  h). \url{https://youtu.be/_V8eKsto3Ug} {[}EN{]}
\item
  \emph{Curso completo de R para Data Science con Tidyverse} {[}Video{]}
  (Juan Gabriel Gomila Salas).
  \url{https://www.youtube.com/watch?v=9NJyhs5PlGQ} {[}ES{]}
  \emph{{[}canal y fecha POR VERIFICAR{]}}
\item
  StatQuest with Josh Starmer. \emph{Linear Regression in R,
  Step-by-Step} {[}Video{]}.
  \url{https://www.youtube.com/watch?v=u1cc1r_Y7M0} {[}EN{]}
\end{itemize}

\subsection*{Nota de verificación}
\addcontentsline{toc}{subsection}{Nota de verificación}

Las referencias se verificaron mediante búsqueda web (título, autores,
editorial, ISBN, DOI y existencia de la URL en índices de búsqueda) el
2026-09-25. No fue posible abrir directamente las páginas desde el
entorno de trabajo, por lo que las \textbf{fechas de publicación de los
videos} y el \textbf{nombre exacto del canal} del video en español
quedan marcados como \emph{POR VERIFICAR}. Revísalos antes de publicar
el curso.
